% Independently written draft from Illusie I, 2.1.2.1, printed pp.23--24.
% Local labels only. Source snippet for bounded root-chapter composition.

\section{Homotopy sheaves of simplicial sheaves}
\label{section-illusie-I-homotopy-sheaves}

\noindent
Let $\mathcal{C}$ be a site. A pointed simplicial sheaf on $\mathcal{C}$
is a simplicial object $X$ of $\Sh(\mathcal{C})$ together with a section
$x : * \to X$. For $n \geq 0$ we write
$\pi_n(X, x)$ for the sheaf associated to
$$
U \longmapsto \pi_n(X(U), x|_U).
$$
For $n = 0$ this is a sheaf of pointed sets, for $n = 1$ a sheaf of
groups, and for $n \geq 2$ a sheaf of abelian groups.

\begin{lemma}
\label{lemma-illusie-I-homotopy-sheaves-colimits-pullback}
Let $n \geq 0$.
\begin{enumerate}
\item If $(X_i, x_i)$ is a filtered diagram of pointed simplicial sheaves
on $\mathcal{C}$, then the canonical map
$$
\colim_i \pi_n(X_i, x_i)
\longrightarrow
\pi_n(\colim_i X_i, \colim_i x_i)
$$
is an isomorphism.
\item If
$f : \Sh(\mathcal{D}) \to \Sh(\mathcal{C})$
is a morphism of topoi, then the canonical map
$$
f^{-1}\pi_n(X, x)
\longrightarrow
\pi_n(f^{-1}X, f^{-1}x)
$$
is an isomorphism.
\item In particular, for a point $p$ of $\Sh(\mathcal{C})$ there is a
canonical isomorphism
$$
\pi_n(X, x)_p = \pi_n(X_p, x_p).
$$
\end{enumerate}
\end{lemma}

\begin{proof}
We first recall exactly which limits and colimits occur in the usual
simplicial-set construction. For a simplicial set $Y$, the standard
fibrant replacement is
$$
\operatorname{Ex}^{\infty}(Y) =
\colim_{r \geq 0} \operatorname{Ex}^r(Y).
$$
In degree $m$ one has
$$
\operatorname{Ex}(Y)_m =
\Mor(\operatorname{Sd}\Delta[m], Y).
$$
The simplicial set $\operatorname{Sd}\Delta[m]$ is finite. Consequently
this mapping set is a finite limit of finitely many of the sets $Y_j$.
Thus every degree of $\operatorname{Ex}^{\infty}(Y)$ is obtained from
$Y$ by iterating finite limits and then taking one filtered colimit.

The simplicial set $\operatorname{Ex}^{\infty}(Y)$ is Kan and
$Y \to \operatorname{Ex}^{\infty}(Y)$ induces isomorphisms on homotopy
groups. If $Z$ is Kan with base point $z$, let $C_n(Z, z)$ be the set of
$n$-simplices all of whose faces are the appropriate degeneracy of $z$.
Let $W_n(Z, z)$ be the set of $(n + 1)$-simplices giving a homotopy between
two elements of $C_n(Z, z)$ relative to the boundary. Both $C_n$ and $W_n$
are finite limits of degrees of $Z$. The two end faces give maps
$$
W_n(Z, z) \rightrightarrows C_n(Z, z),
$$
and their coequalizer is $\pi_n(Z, z)$. Indeed, for a Kan complex the
relation generated by these elementary relative homotopies is already the
usual relative-boundary homotopy relation. We conclude that
$\pi_n(Y, z)$ is constructed functorially using finite limits, colimits,
and filtered colimits only.

Colimits of sheaves are obtained by sheafifying colimits of presheaves,
see Sites, Lemma \ref{sites-lemma-colimit-sheaves}, and sheafification
commutes with finite limits, see Sites, Lemma
\ref{sites-lemma-sheafification-exact}. Hence the same construction made
inside $\Sh(\mathcal{C})$ is the sheaf associated to
$U \mapsto \pi_n(X(U), x|_U)$.

Filtered colimits of sets commute with finite limits by
Categories, Lemma \ref{categories-lemma-directed-commutes}. Colimits
commute with colimits. Applying these facts before sheafification, and
then using the preceding paragraph, proves (1).

The functor $f^{-1}$ commutes with all colimits because it is a left
adjoint, by Categories, Lemma \ref{categories-lemma-adjoint-exact}, and
it commutes with finite limits by the definition of a morphism of topoi,
Sites, Definition \ref{sites-definition-topos}. It therefore commutes
with every operation in the displayed construction of $\pi_n$, which
proves (2). A point is a morphism from the punctual topos; applying (2)
to its inverse-image functor proves (3).
\end{proof}
